% GENERATED FILE. Edit canonical lessons, then run npm run generate:books.
% canonical-source-hash: fnv1a64:9d875a2d3aea638f
% canonical-lessons: ZH-C332-danzi, ZH-C332-gui, ZH-W332-gui-ghost, ZH-C332-naogui, ZH-C332-jiangui, ZH-C332-renjian

\chapter{Ghost Stories}
\label{ch:zh-ghost-stories}
\hlchaptermodality{332}

\begin{chapteropening}
\textbf{By the end of this chapter:} \emph{I can talk about courage, ghosts and a haunted house, curse my luck and talk about the world of the living.}

Complete the last lesson of chapter 332: I can talk about courage, ghosts and a haunted house, curse my luck and talk about the world of the living.
\end{chapteropening}

\section[dǎnzi]{dǎnzi — courage; nerve}

\label{lesson:ZH-C332-danzi}



\begin{quote}
Before the new one: say the Mandarin for a precious treasure, then the Mandarin for to hide away.
\end{quote}

\subsection*{You'll want to know: dǎnzi}
\emph{dǎnzi} — \textquotedblleft{}courage; nerve\textquotedblright{}.

\begin{quote}
\emph{xiǎoshíhou wǒ de dǎnzi hěn xiǎo} — When I was little I was very timid.
\end{quote}

Its characters come later; for now it is a word for the ear.

\begin{cousinweb}
\emph{dǎnzi} is built from \emph{dǎn} 'the gall' and \zh{子}, here a light noun ending, both inherited from Middle Chinese, \emph{tamX} and \emph{tsiX} (Guangyun, in Baxter's transcription), which pairs the first with the liver.
\end{cousinweb}

\begin{grammarlens}[title={What you've built}]
One more word for telling what happened.
\end{grammarlens}

\subsection*{Your turn}
\begin{itemize}
\raggedright
  \item \emph{Say it:} \emph{dǎnzi}
  \item \emph{Say it:} \emph{dǎnzi}, once more
  \item \emph{Say it:} \emph{xiǎoshíhou wǒ de dǎnzi hěn xiǎo} — When I was little I was very timid.
  \item \emph{Recall:} say the Mandarin for a precious treasure, then the Mandarin for to hide away, then say \emph{dǎnzi} again
\end{itemize}

\begin{tcolorbox}[breakable,colback=teal!4,colframe=teal!35!black,title={Before you move on}]
What does dǎnzi mean? (\textquotedblleft{}Courage; nerve\textquotedblright{}.) Say it once more.
\end{tcolorbox}
\section[\texorpdfstring{\zh{鬼} (guǐ)}{鬼 (guǐ)}]{\zh{鬼} (guǐ) — ghost}

\label{lesson:ZH-C332-gui}



\begin{quote}
Before the new one: say the Mandarin for to hide away, then the Mandarin for the nerve to face danger.
\end{quote}

\subsection*{You'll want to know: \zh{鬼}}
\textbf{\zh{鬼}} — \emph{guǐ} — \textquotedblleft{}ghost\textquotedblright{}.

\begin{quote}
\textbf{\zh{那个房子里有鬼}} \emph{nàge fángzi lǐ yǒu guǐ} — There's a ghost in that house.
\end{quote}

The new character, in pieces: \zh{鬼} — nine strokes: a short stroke falling left over a box with a bar inside; from the box two long strokes, one falling to the left, one bending right along the foot and hooking up; and on the right a small turning stroke and a dot. The next lesson puts it in your hand.

\begin{cousinweb}
\textbf{\zh{鬼}} (\emph{guǐ}) is inherited from Middle Chinese \emph{kjw+jX} (Guangyun, in Baxter's transcription), which explains the word by one of like sound, 'to return home'.
\end{cousinweb}

\begin{grammarlens}[title={What you've built}]
One more word for telling what happened.
\end{grammarlens}

\subsection*{Your turn}
\begin{itemize}
\raggedright
  \item \emph{Say it:} \emph{guǐ}
  \item \emph{Say it:} \emph{guǐ}, once more
  \item \emph{Say it:} \emph{nàge fángzi lǐ yǒu guǐ} — There's a ghost in that house.
  \item \emph{Recall:} say the Mandarin for to hide away, then the Mandarin for the nerve to face danger, then say \emph{guǐ} again
\end{itemize}

\begin{tcolorbox}[breakable,colback=teal!4,colframe=teal!35!black,title={Before you move on}]
What does \zh{鬼} mean? (\textquotedblleft{}Ghost\textquotedblright{}.) Say it once more.
\end{tcolorbox}
\section[\texorpdfstring{\zh{鬼} (guǐ)}{鬼 (guǐ)}]{\zh{鬼} — a box with long legs}

\label{lesson:ZH-W332-gui-ghost}



\begin{quote}
The word first, then the character inside it.

\begin{itemize}
\raggedright
  \item \emph{Recall:} say \emph{guǐ} — \textbf{R1}, one lesson back, before it decays
\end{itemize}

\begin{quote}
\zh{口} \zh{儿}
\end{quote}

\emph{Read it:} each one aloud, with nothing said first — \emph{kǒu}, \emph{ér}
\end{quote}

\subsection*{Script — a box with long legs}
\hlblockfigure{\detokenize{figures/ZH-W332-gui-ghost-filmstrip.pdf}}{How \zh{鬼} is written, stroke by stroke}

\begin{quote}
\zh{鬼}
\end{quote}

\textbf{Three pieces: the box, then the legs, then the small piece on the right.}

\textbf{The box, five strokes.} A short stroke falling to the left on top. Then a short vertical on the left, a bar that runs right and turns down, a bar across the middle, and the bar that closes the foot.

\textbf{The legs, two strokes.} A long stroke that starts at the top of the box, runs down through its middle and falls away to the left. Then a stroke that runs down, bends right along the foot and hooks up at the end.

\textbf{The small piece, two strokes,} inside the bend on the right: a short stroke that falls to the left and turns right, then a dot.

5 + 2 + 2 = nine.

\subsection*{Your turn}
\begin{itemize}
\raggedright
  \item \emph{Write it:} \zh{鬼}, all nine, with the model in front of you
  \item \emph{Cover:} write it again from memory, and count the pieces, not the strokes
\end{itemize}

\begin{tcolorbox}[breakable,colback=teal!4,colframe=teal!35!black,title={Before you move on}]
How many strokes in \zh{鬼}? (\textbf{Nine}.) How many pieces? (\textbf{Three: the box, the legs, the small piece}.) Which stroke hooks up? (\textbf{The seventh, along the foot}.)

Source: \href{https://raw.githubusercontent.com/chanind/hanzi-writer-data/68d10a4b21150cae5e1ebbd223eed289cf32d90c/data/\%E9\%AC\%BC.json}{Hanzi Writer Data \zh{鬼}, pinned stroke paths and medians}.
\end{tcolorbox}
\section[nàoguǐ]{nàoguǐ — to be haunted}

\label{lesson:ZH-C332-naogui}



\begin{quote}
Before the new one: say the Mandarin for the nerve to face danger, then the Mandarin for a ghost in a story.
\end{quote}

\subsection*{You'll want to know: nàoguǐ}
\emph{nàoguǐ} — \textquotedblleft{}to be haunted\textquotedblright{}.

\begin{quote}
\emph{dàjiā dōu shuō nàge fángzi nàoguǐ} — Everyone says that house is haunted.
\end{quote}

Its characters come later; for now it is a word for the ear.

\begin{cousinweb}
\emph{nàoguǐ} is built from \emph{nào} 'to make a din' and \zh{鬼} 'a ghost', both inherited from Middle Chinese, \emph{nraewH} and \emph{kjw+jX} (Guangyun, in Baxter's transcription), which glosses the first 'not quiet' and 'to disturb'.
\end{cousinweb}

\begin{grammarlens}[title={What you've built}]
One more word for telling what happened.
\end{grammarlens}

\subsection*{Your turn}
\begin{itemize}
\raggedright
  \item \emph{Say it:} \emph{nàoguǐ}
  \item \emph{Say it:} \emph{nàoguǐ}, once more
  \item \emph{Say it:} \emph{dàjiā dōu shuō nàge fángzi nàoguǐ} — Everyone says that house is haunted.
  \item \emph{Recall:} say the Mandarin for the nerve to face danger, then the Mandarin for a ghost in a story, then say \emph{nàoguǐ} again
\end{itemize}

\begin{tcolorbox}[breakable,colback=teal!4,colframe=teal!35!black,title={Before you move on}]
What does nàoguǐ mean? (\textquotedblleft{}To be haunted\textquotedblright{}.) Say it once more.
\end{tcolorbox}
\section[\texorpdfstring{\zh{见鬼} (jiànguǐ)}{见鬼 (jiànguǐ)}]{\zh{见鬼} (jiànguǐ) — how absurd; curse it}

\label{lesson:ZH-C332-jiangui}



\begin{quote}
Before the new one: say the Mandarin for a ghost in a story, then the Mandarin for to be haunted by a ghost.
\end{quote}

\subsection*{You'll want to know: \zh{见鬼}}
\textbf{\zh{见鬼}} — \emph{jiànguǐ} — \textquotedblleft{}how absurd; curse it\textquotedblright{}.

\begin{quote}
\textbf{\zh{见鬼了我的手机不见了}} \emph{jiànguǐ le wǒ de shǒujī bú jiàn le} — What on earth? My phone's gone!
\end{quote}

\begin{cousinweb}
\textbf{\zh{见鬼}} (\emph{jiànguǐ}) is built from \zh{见} 'to see' and \zh{鬼} 'a ghost', both inherited from Middle Chinese, \emph{kenH} and \emph{kjw+jX} (Guangyun, in Baxter's transcription), which glosses the first 'to look'. The exclamation is not given there.
\end{cousinweb}

\begin{grammarlens}[title={What you've built}]
One more word for telling what happened.
\end{grammarlens}

\subsection*{Your turn}
\begin{itemize}
\raggedright
  \item \emph{Say it:} \emph{jiànguǐ}
  \item \emph{Say it:} \emph{jiànguǐ}, once more
  \item \emph{Say it:} \emph{jiànguǐ le wǒ de shǒujī bú jiàn le} — What on earth? My phone's gone!
  \item \emph{Recall:} say the Mandarin for a ghost in a story, then the Mandarin for to be haunted by a ghost, then say \emph{jiànguǐ} again
\end{itemize}

\begin{tcolorbox}[breakable,colback=teal!4,colframe=teal!35!black,title={Before you move on}]
What does \zh{见鬼} mean? (\textquotedblleft{}How absurd; curse it\textquotedblright{}.) Say it once more.
\end{tcolorbox}
\section[\texorpdfstring{\zh{人间} (rénjiān)}{人间 (rénjiān)}]{\zh{人间} (rénjiān) — the human world}

\label{lesson:ZH-C332-renjian}



\begin{quote}
Before the new one: say the Mandarin for to be haunted by a ghost, then the Mandarin for how absurd.
\end{quote}

\subsection*{You'll want to know: \zh{人间}}
\textbf{\zh{人间}} — \emph{rénjiān} — \textquotedblleft{}the human world\textquotedblright{}.

\begin{quote}
\textbf{\zh{传说七月的时候鬼会来人间}} \emph{chuánshuō qīyuè de shíhou guǐ huì lái rénjiān} — Legend says that in the seventh month ghosts come into the world of the living.
\end{quote}

\begin{cousinweb}
\textbf{\zh{人间}} (\emph{rénjiān}) is built from \zh{人} 'people' and \zh{间} 'between', both inherited from Middle Chinese, \emph{nyin} and \emph{kean} (Guangyun, in Baxter's transcription); the dictionary lists the second under an older form of the character and glosses it 'a gap'.
\end{cousinweb}

\begin{grammarlens}[title={What you've built}]
One more word for telling what happened.
\end{grammarlens}

\subsection*{Your turn}
\begin{itemize}
\raggedright
  \item \emph{Say it:} \emph{rénjiān}
  \item \emph{Say it:} \emph{rénjiān}, once more
  \item \emph{Say it:} \emph{chuánshuō qīyuè de shíhou guǐ huì lái rénjiān} — Legend says that in the seventh month ghosts come into the world of the living.
  \item \emph{Recall:} say the Mandarin for to be haunted by a ghost, then the Mandarin for how absurd, then say \emph{rénjiān} again
\end{itemize}

\begin{tcolorbox}[breakable,colback=teal!4,colframe=teal!35!black,title={Before you move on}]
What does \zh{人间} mean? (\textquotedblleft{}The human world\textquotedblright{}.) Say it once more.
\end{tcolorbox}
